\section{Discussion}
\label{sec:discussion}

\emph{Assumptions.} The operational equivalences used here follow from the quantum description of the decay: a parent of definite spin $j$, rotation covariance, and hence angular distributions of degree at most $2j$. Under reading I they are properties of the data and can be checked on the data, for instance by the absence of higher harmonics; but an equivalence is a statement about all preparations, and only those of the fragment can be examined. This is the tomographic-completeness assumption common to noncontextuality tests \cite{mazurek2016,pusey2018}. The theorems also assume full angular acceptance; selections that depend on the recorded directions break covariance and must be removed by unfolding, as in the entanglement analyses.

\emph{What is excluded.} A failure of noncontextuality is not a loophole-free exclusion of local hidden variables: the recorded-outcome model always exists. What is excluded is every hidden-variable account in which the response to a set of outcomes depends only on the effect that the set defines, equivalently every local model with polynomial response functions of degree $2j$ [Eq.~\eqref{eq:pairpoly}]. Whether this is the appropriate notion of classicality for collider data is a matter of interpretation. In its favour: it is the notion equivalent to simplex embedding and to nonnegative quasiprobability representations \cite{schmid2021,spekkens2008}, it reduces to separability when all local measurements are available \cite{plavala2024}, and for decays it gives conditions on measured coefficients with thresholds that are neither trivial nor identical to those of entanglement.

\emph{Relation to entanglement measurements.} For decay data, contextuality implies entanglement and is strictly stronger unless both analysing powers are maximal. It is tested on raw angular coefficients and needs large analysing powers on both sides: $|\kappa|>1/3$ for spin $1/2$, kernels outside $\cR$ for spin $1$. The three examples of Sec.~\ref{sec:apps} show the three possibilities: for top-quark pairs with two charged leptons the tests coincide; for $J/\psi\to\Lambda\bar\Lambda$ the state is entangled everywhere and the data are contextual in a window; for pairs analysed through $b$ quarks, or through a $Z\to\ell\ell$ decay of a spin-$1$ partner, the data are noncontextual whatever the state.

\emph{Open problems.}
(1) Spin-$1$ pairs outside the isotropic family, in particular $H\to WW^*$ and $H\to ZZ^*$. By Proposition~\ref{prop:pairpoly} the question is whether a polynomial of bidegree $(2,2)$ lies in the cone generated by products of nonnegative quadratic polynomials on two spheres; Lemma~\ref{lem:sos} describes one factor, and the linear programmes of Appendix~\ref{app:numerics} decide instances.
(2) The exact region of classical kernels for $j\ge\tfrac32$; Theorem~\ref{thm:spinj} gives its extent in the vector direction only.
(3) Whether the extremal kernel $D^\star$ has nonzero multipole ratios for all $j$, which would make the bound of Theorem~\ref{thm:spinj} a maximum for all $j$.
(4) Exactly degenerate kernels (Remark~\ref{rem:deg}) and, with them, the robustness of the criteria to imperfect equivalences and finite statistics; noise-robust inequalities follow from the duals of the linear programmes \cite{selby2024lp,kunjwal2019}, and the witnesses of Corollary~\ref{cor:D} and Eq.~\eqref{eq:dipole} are the simplest examples.
(5) Decays with recorded energies and cascades such as $\Xi\to\Lambda\pi\to p\pi\pi$, where the polarimeter has more than one stage.
